% =====================================================================
%  chunks/rev_ch02.tex —— 第 2 章 读书笔记（tr-wu 撰）
% =====================================================================
\chapterreview{%
第 2 章只占原书 p.15--40 不足三十页，却是全书的语法书：它从「拓扑空间」一路推到「纤维化的长正合列」，为第 3 章的障碍论与第 4 章的流形分类备好全部语言。作者在此不追求计算，几乎所有结论都停在「存在某个同构」的层面，真正的计算留给同调论；因此本章的正确读法是把它当字典用，而不是当习题集刷。

脉络极清楚：§1 用开集重述连续、紧、Hausdorff、乘积拓扑与楔和，并把映射柱 (1.2) 与拓扑流形的定义 1.1 一次给足；§2 把连续形变形式化为同伦，给出同伦等价、形变收缩与可缩；§3 引入覆叠同伦性质，得到纤维化定义 3.1 与同伦联络，并证明同一道路连通分支上的纤维彼此同伦等价；§4 才是重头戏，同伦群的定义 4.1、群环的公式 (4.1)(4.2) 与相对同伦群的定义 4.2，随即导出正合性 (4.4)、(4.5) 的正合同伦序列，以及 (4.7) 的纤维化长正合列。

核心结论应当记住编号：(4.8) 说覆叠空间的投影在 $n>1$ 时诱导同伦群同构，(4.9) 说 $\pi_n(Y)\cong\pi_{n-1}(\Omega Y)$，(4.6) 把相对同伦群换回绝对同伦群。这三条合起来，就是「用覆叠与纤维化反推同伦群」的流水线，后面算 $S^n$、$\RP^n$、$\CP^n$ 与经典群全靠它。(4.22) 断言除 $S^2$、$\RP^2$ 外曲面高阶同伦群全平凡，(4.23) 定义 Eilenberg--MacLane 空间 $K(\pi,n)$，(4.25) 给出 $\pi_3(S^2)\cong\Z$，这是 Hopf 纤维化（书中未编号；导出的一组同构即 (4.24)）的直接收益。

辨析有三处坑。其一，同伦等价远弱于同胚：$\R^n$ 可缩，单位球与点是同一同伦型，切勿把同伦型当成「形状」。其二，纤维化与覆叠不是一回事：覆叠要求离散纤维与唯一提升，纤维化的纤维可任意，只要求同伦提升；(4.8) 这类同构都暗中用到了纤维离散。其三，交换性有边界：绝对同伦群 $\pi_n$（$n>1$）交换，相对同伦群 $\pi_n(X,A)$ 要到 $n>2$ 才交换，而 $n=1$ 时 $\pi_1$ 的共轭作用恰是「基点不可忽略」的根源。

译注时踩到几处原书笔误，读者宜警惕：p.24 把 $S^n$ 写成 $\sum(x^\alpha)^2\le 1$（应为 $=1$，否则是球体）；p.29 相对同伦群定义中印作 $f(S^n)\subset A$，按定义域应为 $f(S^{n-1})\subset A$；p.21 关于覆叠空间「所有子集都是开集」的英文原句有脱漏。这些一律照原书排印，另加批注说明。

总评：本章的写法是「定义—正合列—八个例子」三段式，例子从 $\RP^2$、环面、曲面一直排到 Hopf 纤维化、纽结补集与 $\CP^\infty$，几乎覆盖第 3、4 章的全部常用模型。缺点是存在性定理多只声明不证明（万有覆叠的存在、$K(\pi,n)$ 的 CW 实现），初学者宜配 Hatcher 前两章对照；但作为以语言准备为目的的综述，其取舍极经济：后面要用的这里全给，不用的绝不铺陈。%
}

% CHUNK rev_ch02 DONE
